\documentclass[11pt,a4paper]{article}
\usepackage[utf8]{inputenc}
\usepackage[T1]{fontenc}
\usepackage[french]{babel}
\usepackage{amsmath, amssymb}
\usepackage{geometry}
\geometry{margin=2.5cm}
\usepackage{xcolor}
\usepackage{listingsutf8}
\usepackage[most]{tcolorbox}
\usepackage{prftree}
\usepackage{hyperref}
\usepackage{newunicodechar}

\newunicodechar{ι}{\ensuremath{\iota}}
\newunicodechar{₁}{\ensuremath{_1}}
\newunicodechar{₂}{\ensuremath{_2}}

\newtcbtheorem{remarque}{Remarque}{
  colback=blue!5!white,
  colframe=blue!75!black,
  fonttitle=\bfseries,
  breakable,
  enhanced
}{rem}

\newtcbtheorem{exercice}{Exercice}{
  colback=green!5!white,
  colframe=green!40!black,
  fonttitle=\bfseries,
  breakable,
  enhanced
}{exo}

\definecolor{backcolour}{rgb}{0.95,0.95,0.92}
\definecolor{codegreen}{rgb}{0,0.6,0}
\lstset{
  language=ML,
  inputencoding=utf8,
  backgroundcolor=\color{backcolour},
  commentstyle=\color{codegreen},
  keywordstyle=\color{blue}\bfseries,
  basicstyle=\ttfamily\small,
  breaklines=true,
  frame=single,
  rulecolor=\color{lightgray},
  showstringspaces=false,
  extendedchars=true,
  literate={ι}{{$\iota$}}1 {₁}{{$_1$}}1 {₂}{{$_2$}}1 {é}{{\'e}}1 {è}{{\`e}}1 {à}{{\`a}}1 {ù}{{\`u}}1 {ê}{{\^e}}1 {î}{{\^i}}1 {ô}{{\^o}}1 {û}{{\^u}}1 {ë}{{\"e}}1 {ï}{{\"i}}1 {ö}{{\"o}}1 {ü}{{\"u}}1
}

\makeatletter
\newcommand{\subtitle}[1]{%
  \gdef\@subtitle{#1}%
}
\let\oldmaketitle\maketitle
\renewcommand{\maketitle}{%
  \let\oldtitle\@title
  \renewcommand{\@title}{\oldtitle\\[1em]\Large\normalfont\@subtitle}%
  \oldmaketitle
}
\makeatother

\title{LEÇON 4 : Proof-theoretic semantics, vérificationnisme et pragmatisme}
\subtitle{Éliminateurs à la main, propriétés de l'égalité, et \texttt{Compute}}
\date{07/04/2026}
\author{Pablo Donato}

\begin{document}

\maketitle

Nous avons vu aujourd'hui que les connecteurs logiques sont définis dans Rocq comme des \textbf{types inductifs}, dont les constructeurs sont exactement les règles d'introduction de la déduction naturelle. La tactique \texttt{destruct} génère automatiquement un \texttt{match~\dots~with~\dots~end} pour les éliminer. Dans ces exercices, vous écrirez ce \texttt{match} à la main pour voir ce qui se passe sous le capot. Vous décortiquerez ensuite les propriétés fondamentales de l'égalité, dont le comportement est lui aussi justifié par la sémantique vérificationniste.

\begin{lstlisting}
Section Lesson4.
Variables A B C : Prop.
\end{lstlisting}

\section{Éliminateurs à la main (\texttt{match~\dots~with})}

Rappel : le type \texttt{A /\textbackslash{} B} a un seul constructeur \texttt{conj~:~A~->~B~->~A~/\textbackslash{}~B}. Pour utiliser une hypothèse \texttt{H~:~A~/\textbackslash{}~B}, on ``pattern-matche'' ou ``filtre par motif'' :

\begin{lstlisting}
match H with
| conj a b => (* a : A, b : B disponibles ici *)
end
\end{lstlisting}

Le type \texttt{A \textbackslash{}/ B} a deux constructeurs : \texttt{or\_introl~:~A~->~A~\textbackslash{}/~B} et \texttt{or\_intror~:~B~->~A~\textbackslash{}/~B}. Le filtrage génère deux branches :

\begin{lstlisting}
match H with
| or_introl a => ...
| or_intror b => ...
end
\end{lstlisting}

Voici un exemple complet pour illustrer la démarche : on extrait les composantes avec \texttt{match}, puis on reconstruit avec \texttt{conj}.

\begin{lstlisting}
(* Exemple : symetrie de la conjonction *)
Lemma and_sym : A /\ B -> B /\ A.
Proof.
  intro H.
  exact (match H with
         | conj a b => conj b a
         end).
Qed.
\end{lstlisting}

\begin{exercice}{Extraction de la deuxième composante}{}
  Sans utiliser \texttt{destruct}, prouvez que $A \land B \Rightarrow B$ en écrivant un \texttt{match} explicite.

\begin{lstlisting}
Lemma and_elim_2 : A /\ B -> B.
Proof.
  (* A COMPLETER *)
Admitted.
\end{lstlisting}
\end{exercice}

\begin{exercice}{Symétrie de la disjonction}{}
  Sans utiliser \texttt{destruct}, prouvez $A \lor B \Rightarrow B \lor A$.
  Votre \texttt{match} doit avoir deux branches — l'une utilisant \texttt{or\_intror}, l'autre \texttt{or\_introl}.

\begin{lstlisting}
Lemma or_sym : A \/ B -> B \/ A.
Proof.
  intro H.
  (* A COMPLETER *)
Admitted.
\end{lstlisting}
\end{exercice}

\begin{exercice}{De la conjonction vers la disjonction}{}
  Utilisez \texttt{match} pour extraire $A$ de l'hypothèse $A \land B$, puis injectez-le dans $A \lor C$ avec \texttt{or\_introl}.

\begin{lstlisting}
Lemma and_to_or : A /\ B -> A \/ C.
Proof.
  (* A COMPLETER *)
Admitted.
\end{lstlisting}
\end{exercice}

\begin{remarque}{Correspondance avec \texttt{destruct}}{}
  À tout moment, vous pouvez utiliser \texttt{Print} pour voir le terme de preuve réel que Rocq a construit :
\begin{lstlisting}
Print and_sym.
\end{lstlisting}
  Vous constaterez que \texttt{destruct H as [a b]} génère exactement le même \texttt{match} que vous venez d'écrire à la main. La tactique n'est qu'un raccourci syntaxique.
\end{remarque}

\section{Propriétés de l'égalité — $\lambda$-termes à la main}

L'éliminateur de \texttt{eq} dérivé automatiquement par Rocq est \texttt{eq\_ind} :
\begin{lstlisting}
eq_ind : forall (A : Type) (x : A) (P : A -> Prop),
         P x -> forall y : A, x = y -> P y
\end{lstlisting}
C'est le \textbf{principe de Leibniz} : si $P(x)$ est vrai et $x = y$, alors $P(y)$ est vrai.
La tactique \texttt{rewrite~H} l'applique silencieusement.

\textbf{Stratégie :} pour prouver $Q(y)$ sachant \texttt{H~:~x~=~y}, choisir $P := \lambda z.\; Q(z)$ de sorte que $P(x) = Q(x)$, puis écrire \texttt{eq\_ind~x~P~(preuve de Q x)~y~H}.

L'exemple de la symétrie (vu dans \texttt{cours.v}) illustre la démarche :
\begin{lstlisting}
(* Exemple : symetrie.  P := fun z => z = x, donc P x = (x = x). *)
Definition eq_sym_terme : forall (x y : T), x = y -> y = x :=
  fun x y H => eq_ind x (fun z => z = x) eq_refl y H.
\end{lstlisting}

\begin{lstlisting}
Variable T : Type.
\end{lstlisting}

\begin{exercice}{Transitivité — $\lambda$-terme pur}{}
  Pour prouver $x = z$ sachant $H_1 : x = y$ et $H_2 : y = z$, choisir $P := \lambda w.\; x = w$.
  Alors $P(y) = (x = y)$, prouvé par $H_1$.
  Appliquer \texttt{eq\_ind} à $H_2 : y = z$ pour obtenir $P(z) = (x = z)$.

\begin{lstlisting}
Definition eq_trans : forall (x y z : T), x = y -> y = z -> x = z :=
  todo. (* A COMPLETER *)
\end{lstlisting}
\end{exercice}

\begin{exercice}{Congruence — $\lambda$-terme pur}{}
  Si $x = y$, alors $f(x) = f(y)$ pour toute $f : T \to T$.
  Choisir $P := \lambda z.\; f(x) = f(z)$, de sorte que $P(x) = (f(x) = f(x))$,
  prouvé par \texttt{eq\_refl}.

\begin{lstlisting}
Definition f_equal : forall (f : T -> T) (x y : T), x = y -> f x = f y :=
  todo. (* A COMPLETER *)
\end{lstlisting}
\end{exercice}

\section{Expérimenter avec \texttt{Compute}}

La commande \texttt{Compute} évalue une expression jusqu'à sa forme normale. Elle illustre l'\textbf{élimination des détours} (\textit{cut elimination}) : un constructeur immédiatement suivi de son éliminateur se réduit en une étape.

\begin{lstlisting}
Variable (a : A) (b : B) (f : A -> C) (g : B -> C).
\end{lstlisting}

Exécutez les commandes suivantes et observez les résultats :

\begin{lstlisting}
(* Detour fonctionnel (beta) *)
Compute (fun x : A => x) a.

(* Detour conjonctif *)
Compute (match conj a b with conj x y => x end).

(* Detour disjonctif *)
Compute (match or_introl a : A \/ B with
         | or_introl x => f x
         | or_intror y => g y
         end).
\end{lstlisting}

\begin{exercice}{Observer la réduction}{}
  \begin{enumerate}
    \item Que retourne chacune des trois commandes \texttt{Compute} ci-dessus ?
    \item Que se passe-t-il si vous remplacez \texttt{or\_introl~a} par \texttt{or\_intror~b} dans le troisième exemple ?
    \item \textit{(Bonus)} Construisez un exemple de ``double détour'' — deux introductions suivies de deux éliminations — et vérifiez avec \texttt{Compute} que la normalisation donne directement la valeur finale.
  \end{enumerate}
\end{exercice}

\begin{lstlisting}
End Lesson4.
\end{lstlisting}

\end{document}
